\subsection{Tail-Transverse Influence}
\label{app:tail-transverse-influence}
\label{sec:targeted-interventions}

Fix the predictor, coordinates, calibration, slice rules and finite audit weights $\pi_g\ge0$, $\sum_g\pi_g=1$. An action specifies a positive-semidefinite path $J_g(t,a)=J_g+tH_{g,a}+o(t)$ as $t\downarrow0$, with symmetric $H_{g,a}$ and an operator-norm remainder. Allocation, normalization and retained pooling belong in this path; audit weights stay fixed.

\paragraph{One differential for exact and ridge targets.}
Write $J_g=\begin{psmallmatrix}A_g&B_g\\B_g^\top&C_g\end{psmallmatrix}$. Choose a fixed $\rho\ge0$ with $C_g+\rho I\succ0$: $\rho=0$ gives exact Schur information and $\rho=\lambda_{\rid}$ gives the ridge target. Set $Q_g=(C_g+\rho I)^{-1}$, $S_\rho(J_g)=A_g-B_gQ_gB_g^\top$, and let $\lambda_g$ be its smallest eigenvalue. Since $DQ_g[H_{NN}]=-Q_gH_{NN}Q_g$,
\begin{equation}
\begin{aligned}
DS_\rho(J_g)[H]={}&H_{TT}-H_{TN}Q_gB_g^\top-B_gQ_gH_{TN}^\top\\
&+B_gQ_gH_{NN}Q_gB_g^\top,\\
p_g(u)&=\begin{pmatrix}u\\-Q_gB_g^\top u\end{pmatrix},
\qquad u^\top DS_\rho(J_g)[H]u=p_g(u)^\top Hp_g(u).
\end{aligned}
\label{eq:influence-schur-differential}
\end{equation}
For a simple minimum eigenvalue with unit eigenvector $u_g$, differentiating the eigenvalue equation therefore gives
\begin{equation}
h_{g,a}:=D_a\lambda_g=p_g^\top H_{g,a}p_g,
\qquad p_g=p_g(u_g).
\label{eq:directional-schur-derivative}
\end{equation}
For allocation rate $\eta_g\ge0$, additive and normalized contributions have different derivatives:
\begin{equation}
\begin{aligned}
H_{g,a}=\eta_gs_as_a^\top
&\ \Longrightarrow\ h_{g,a}=\eta_g(p_g^\top s_a)^2,\\
J_g(t,a)=\frac{J_g+t\eta_gM_{g,a}}{1+t\eta_g}
&\ \Longrightarrow\ H_{g,a}=\eta_g(M_{g,a}-J_g),\\
&\hspace{13mm}h_{g,a}=\eta_g\!\left[p_g^\top M_{g,a}p_g-\lambda_g
+\rho\lVert Q_gB_g^\top u_g\rVert^2\right].
\end{aligned}
\label{eq:influence-ridge-normalized-gain}
\end{equation}
Indeed $Q_gC_gQ_g=Q_g-\rho Q_g^2$, so $p_g^\top J_gp_g=\lambda_g-\rho\lVert Q_gB_g^\top u_g\rVert^2$; the correction vanishes at $\rho=0$. Rank-one gain is the squared efficient residual $[u_g^\top(s_{T,a}-B_gQ_gs_{N,a})]^2$. Replacement instead uses the allocated new-minus-removed contribution. For one observation appended to group $g$, its indicator and original sample-count scaling enter $\eta_g$; they cannot be omitted when claiming a normalized-augmentation derivative.

\paragraph{Tail sensitivity, including nonsmooth cases.}
Fix $\varepsilon>0$ and $0<\alpha<1$, and put $L_g=(\lambda_g+\varepsilon)^{-1}$. Appendix~\ref{app:theory-certificate}'s finite CVaR dual uses
$\mathcal W_{\alpha,\pi}=\{\omega\ge0:\sum_g\omega_g=1,\ \omega_g\le\pi_g/(1-\alpha)\}$.
Let $\Omega_\alpha^\star(L)$ maximize $\sum_g\omega_gL_g$ over this set. If the maximizer $\omega_\alpha$ is unique and eigenvalues are simple,
\begin{equation}
G_\alpha(a):=-D_a\CVaR_{\alpha,\pi}(L)
=\sum_g\omega_{\alpha,g}\frac{h_{g,a}}{(\lambda_g+\varepsilon)^2},
\qquad D_a I_{\alpha,\pi}=I_{\alpha,\pi}^{\,2}G_\alpha(a).
\label{eq:tail-information-gain}
\end{equation}
This follows from $D_aL_g=-h_{g,a}/(\lambda_g+\varepsilon)^2$ and differentiation of the unique active linear functional. The dual weights include fractional threshold mass and the upper-tail normalization.

For a repeated minimum eigenvalue with eigenspace $\mathcal E_g$, symmetric eigenvalue perturbation gives the first formula below. At a CVaR tie, the right derivative of a maximum is the maximum of its active directional derivatives; negation gives the second:
\begin{equation}
\begin{aligned}
h_{g,a}&=\min_{\substack{u\in\mathcal E_g\\\lVert u\rVert=1}}
p_g(u)^\top H_{g,a}p_g(u),\\
G_\alpha^+(a):=-D_a^+\CVaR_{\alpha,\pi}(L)
&=\min_{\omega\in\Omega_\alpha^\star(L)}
\sum_g\omega_g\frac{h_{g,a}}{(\lambda_g+\varepsilon)^2}.
\end{aligned}
\label{eq:influence-tied-cvar-gain}
\end{equation}
These formulas cover either or both sources of nondifferentiability; the reciprocal-certificate identity holds with right derivatives. An arbitrary fixed maximizing tail vector is only a score unless it attains this direction-dependent minimum.

For assigned label $y_a$, tail-transverse influence multiplies the directional information gain by current prediction loss:
\begin{equation}
\TTI(a)=G_\alpha(a)\loss_w(x_a^{\mathrm{cf}},y_a),
\qquad \loss_w(x,y)=-\log p_w(y\mid x),
\label{eq:tail-transverse-influence}
\end{equation}
using $G_\alpha^+$ at nonsmooth points. This ranks actions at the fixed predictor; finite augmentation, refitting and changed normalization require a new audit. The implemented residual-square surrogate $\widetilde G_\alpha$ below omits allocation/baseline terms and is distinguished from this derivative.
